\chapter{Conclusiones y Recomendaciones}
\label{chap:conclusiones}

Este capítulo sostiene que el trabajo cumple lo que puede demostrar con evidencia verificable a la fecha de corte ---27 de octubre de 2026--- y que declara con la misma precisión lo que todavía depende de la aplicación en la Universidad Tecnológica Indoamérica. Por ello, cada conclusión responde a un objetivo específico, distingue entre resultados verificados y resultados pendientes, y no extiende sus afirmaciones más allá del caso de estudio.

\section{Conclusiones}

\textbf{Respecto del objetivo específico 1 (diagnóstico).} Se determinó con precisión qué se necesita para diagnosticar el tratamiento de datos personales del proceso de admisión y matrícula: diez requisitos de información (D01 a D10), cada uno con su fuente, su instrumento y su fundamento normativo (Tabla~\ref{tab:requisitos_diagnostico}), y los instrumentos de recolección correspondientes (Anexos~\ref{anx:lista_verificacion} a~\ref{anx:consentimiento} y~\ref{anx:competencia}). El análisis mostró además que los requisitos no son equivalentes: la autorización institucional (D01) condiciona a todos los demás, y el inventario de tratamientos y el volumen de aspirantes menores de edad (D03 y D04) determinan si el proceso es de gran escala y, con ello, si la evaluación de impacto y el delegado son obligatorios. La línea base de cumplimiento de la UTI ---nivel de madurez, requisitos cumplidos y RAT--- no puede afirmarse todavía, porque depende de la autorización formal, que no se había emitido a la fecha de corte.

\textbf{Respecto del objetivo específico 2 (universidades).} El análisis de fuentes oficiales permitió cuantificar la brecha que la tesis atiende: en el ciclo 2024--2025 funcionaron 16\,161 instituciones educativas, y el sistema de educación superior cuenta con 64 universidades y escuelas politécnicas y 216 institutos superiores, todas responsables del tratamiento según la LOPDP; frente a ello, los delegados de protección de datos registrados en abril de 2025 equivalían como máximo al 7,0\,\% de esas instituciones, y no existe una medición oficial del porcentaje de universidades que ha implementado la ley (Sección~\ref{sec:cifras_sector}). Se verificó, además, que el plazo de notificación a la autoridad es un término de cinco días y el de comunicación a los titulares, de tres días, mientras que el de 72 horas corresponde a incidentes de ciberseguridad de entidades públicas y operadores de infraestructura crítica (Tabla~\ref{tab:plazos_vulneraciones}); que Ecuador se encuentra en una etapa de consolidación regulatoria con aplicación incipiente, y que su régimen es, junto con el español, de los más exigentes para las instituciones educativas entre los seis comparados. Estos hallazgos se tradujeron en siete decisiones de diseño concretas (Tabla~\ref{tab:insumos_oe2}), entre ellas la fase de preparación y su equivalencia con ISO/IEC 27035-1:2023 y NIST SP 800-61r3.

\textbf{Respecto del objetivo específico 3 (implementación).} La metodología quedó lista para implementarse en el proceso de admisión y matrícula: se construyeron trece escenarios candidatos que cubren las cuatro dimensiones de análisis y los cuatro tipos de vulnerabilidad exigidos por la guía de la SPDP, y cada uno quedó vinculado con su requisito normativo, su control de referencia en ISO/IEC 27002:2022, un rol responsable propuesto, la regla que fija su plazo y un indicador de seguimiento (Tabla~\ref{tab:trazabilidad}). Las matrices de riesgo inherente y residual y el plan de tratamiento definitivo dependen de la confirmación de los escenarios en el diagnóstico y de la estimación del panel, que no se habían realizado a la fecha de corte; por ello no se reporta todavía la distribución de niveles.

\textbf{Respecto del objetivo específico 4 (evaluación).} La evaluación se diseñó con doce indicadores con metas fijadas antes de conocer los resultados (Tabla~\ref{tab:indicadores_validacion}) y dos escenarios de validación: la aplicación real (V1) y un simulacro de vulneración (V2). Los tres indicadores técnicos cumplen su meta: los once requisitos funcionales tienen evidencia en el diseño, las 1\,057 fórmulas de la plantilla se evalúan sin errores y los seis casos de prueba se superan, incluidos los que impiden que un daño muy grave para el titular quede oculto tras una probabilidad baja. Los nueve indicadores restantes ---justificación y trazabilidad de las estimaciones, consenso Delphi, calibración y competencia del panel, juicio de expertos y tiempos del simulacro--- no tienen resultado a la fecha de corte, porque dependen de la aplicación en la institución. En consecuencia, este trabajo no afirma todavía que la metodología sea válida; afirma que está completamente especificada, que su herramienta es correcta y que los criterios con los que se juzgará su validez son explícitos y verificables.

\textbf{Limitaciones.} Los resultados se circunscriben al proceso de admisión y matrícula de la UTI, sede Quito, y no se generalizan a otras instituciones ni a otros procesos. Las cifras sectoriales proceden en parte de medios que citan a la fuente oficial y llevan su fecha de corte; varias no se publican y quedan señaladas como pendientes. Las estimaciones dependerán del juicio de expertos, con la incertidumbre propia de ese método, mitigada mediante la calibración y el coeficiente de competencia. Por último, la normativa de la SPDP está en evolución, de modo que la metodología refleja la vigente al 27 de octubre de 2026.

\section{Recomendaciones}

Las recomendaciones se derivan de los hallazgos verificados y se dirigen a la institución del caso:

\begin{enumerate}
  \item Formalizar la autorización del estudio y designar un enlace institucional, porque de ello depende la ejecución de todo el diagnóstico (requisito D01).
  \item Determinar, con los datos agregados de admisión, si el proceso alcanza el supuesto de gran escala y, de confirmarse, designar un delegado de protección de datos y realizar la evaluación de impacto antes del siguiente periodo de admisión.
  \item Aprobar un procedimiento de respuesta a vulneraciones con los canales vigentes (SISPDP, organismo de control y CSIRT) y una meta interna de 72 horas, de modo que el término legal de cinco días funcione como límite y no como objetivo, y comprobarlo con un simulacro anual.
  \item Adoptar la plantilla de matriz como registro institucional de riesgos y revisarla al cierre de cada periodo de admisión y ante cambios de sistemas, proveedores o normativa.
  \item Solicitar a la SPDP, mediante acceso a la información pública, las estadísticas de delegados registrados y de vulneraciones notificadas por el sector educativo, para sustituir los datos pendientes por cifras oficiales.
\end{enumerate}

\section{Trabajos Futuros}

Se identifican cuatro líneas de continuidad. Primera, aplicar la metodología a otros procesos con categorías especiales ---becas, bienestar estudiantil y aula virtual--- para comparar resultados dentro de la misma institución. Segunda, medir el grado de implementación de la LOPDP en las 64 universidades del país mediante un instrumento común, lo que llenaría el vacío estadístico identificado en este trabajo. Tercera, migrar las estimaciones ordinales a un modelo cuantitativo, como FAIR con redes bayesianas, cuando la institución acumule un registro histórico de incidentes suficiente. Cuarta, automatizar la plantilla en un sistema de registro con control de versiones y trazabilidad de cada decisión.

En síntesis, el trabajo entrega una metodología completa, una herramienta verificada y un protocolo de aplicación y evaluación con criterios explícitos; lo que queda por demostrar depende de datos que solo puede aportar la institución, y el documento los identifica uno por uno en lugar de suplirlos con supuestos.
